\section{Análisis de ondas Rayleigh}
\label{sec:metodologia}

Una perturbación aplicada en la superficie genera diferentes tipos de ondas, entre ellas ondas de cuerpo y superficiales. Para analizar la propagación de las ondas Rayleigh se establece una geometría de medición en la que la dirección longitudinal une la fuente con los puntos de observación, la dirección vertical es normal a la superficie y la dirección transversal completa el sistema de coordenadas.

El análisis de la dispersión requiere conservar tanto la evolución temporal del movimiento como la posición correspondiente a cada traza. Si $x_m$ representa la posición longitudinal del receptor y $t_n$ el instante de muestreo, el registro de la componente vertical $u_z$ puede organizarse como
\begin{equation}
	U[m,n]=u_z(x_m,t_n).
	\label{eq:registro-xt}
\end{equation}

Cada fila de esta matriz corresponde a una traza temporal adquirida en una posición determinada, mientras que el conjunto constituye un registro espacio--tiempo. La propagación de una componente coherente se manifiesta como una variación sistemática de su posición temporal entre las distintas trazas. Para que esta información pueda interpretarse correctamente deben conservarse tanto la geometría de adquisición como una referencia temporal común.

El análisis espectral de ondas superficiales (SASW) utiliza dos receptores para estimar la diferencia de fase según la frecuencia. Su sencillez instrumental contrasta con las ambigüedades de ciclos enteros y la limitada redundancia espacial para separar contribuciones superpuestas; por ello se repite con distintas separaciones entre receptores \citep{Foti2014}.

En cambio, MASW extiende este principio a múltiples posiciones: la redundancia espacial permite reconocer patrones coherentes y distinguir contribuciones superpuestas. Los registros adquiridos desde la superficie se procesan para determinar las características dispersivas de las ondas Rayleigh y caracterizar el terreno somero \citep{Park1998,Park1999,Foti2018}.

El MASW convencional adquiere simultáneamente mediante un arreglo, con un transductor y un canal por posición. Para reducir ese costo instrumental puede sintetizarse la dimensión espacial mediante adquisiciones sucesivas. \citet{LinAshlock2016} validaron una configuración con receptor fijo y posiciones sucesivas de la fuente, recíproca de la que desplaza el receptor. En este tipo de adquisiciones, cada traza debe vincularse con la posición del receptor y con la excitación que la originó. La repetibilidad de la fuente y una referencia temporal común son, por tanto, necesarias para integrar los registros sucesivos sin confundir variaciones del terreno con diferencias entre impactos.

La aplicación de MASW establece así los principales requerimientos de la adquisición: conservar la información temporal y espacial de las trazas, mantener una referencia temporal entre registros y preservar adecuadamente la amplitud y la fase de las componentes de interés.

Estos requerimientos constituyen el punto de partida para el diseño del sistema de adquisición presentado en la sección siguiente.
